\documentclass[11pt]{article}
\usepackage[T1]{fontenc}
\usepackage{lmodern}
\usepackage[margin=1in]{geometry}
\usepackage{booktabs}
\usepackage{siunitx}
\usepackage{pgfplotstable}
\pgfplotsset{compat=1.18}
\usepackage{hyperref}

\title{A Short Motion Report}
\author{LaTeX for Everyone}
\date{\today}

\begin{document}
\maketitle
\begin{abstract}
This worked report uses four synthetic position measurements to
demonstrate a reproducible table, a plot, and a simple calculation.
The values are teaching data, not experimental evidence.
\end{abstract}

\section{Question and method}
Does the example data approximately follow constant-speed motion?
We use the times and positions in \texttt{measurements.csv}. Both
the table and the plotted points read that same file.

\section{Results}
\begin{table}[htbp]
  \centering
  \caption{Synthetic positions recorded at one-second intervals.}
  \label{tab:measurements}
  \pgfplotstabletypeset[
    col sep=comma,
    columns/time/.style={column name={Time (s)}},
    columns/position/.style={column name={Position (m)},
      fixed,fixed zerofill,precision=1},
    every head row/.style={before row=\toprule,after row=\midrule},
    every last row/.style={after row=\bottomrule}
  ]{measurements.csv}
\end{table}

\begin{figure}[htbp]
  \centering
  \begin{tikzpicture}
    \begin{axis}[width=0.8\linewidth,
        xlabel={Time (s)},ylabel={Position (m)},
        grid=major,legend pos=north west]
      \addplot[only marks,mark=*,black]
        table[x=time,y=position,col sep=comma]{measurements.csv};
      \addlegendentry{Synthetic measurements}
      \addplot[blue,dashed,domain=0:3] {2*x};
      \addlegendentry{Reference line: position = 2 time}
    \end{axis}
  \end{tikzpicture}
  \caption{The points lie near the dashed constant-speed reference line.}
  \label{fig:motion}
\end{figure}

\section{Interpretation and limitations}
From Table~\ref{tab:measurements}, the endpoint estimate is
\[
  v_{\mathrm{avg}}=\frac{6.0-0.0}{3-0}
    =\qty{2.0}{\metre\per\second}.
\]
Figure~\ref{fig:motion} is consistent with that simple reference,
but four synthetic points cannot establish a physical law or
quantify measurement uncertainty.

\section{Conclusion}
The example shows how one data file can drive a report's table
and plot. Real studies need justified sampling, uncertainty analysis,
and evidence appropriate to their research question.
\end{document}
